\section{Conclusion}
\label{sec:conclusion}

SCR, WSCR, ESCR and SDSCR were followed through a week of operating states on a 37-bus benchmark
with three IBR plants, intact and under every single line and transformer outage, and each
change was split into its voltage, interaction and driving-point-impedance parts. Wherever
commitment varies, the four metrics agree on when each site is weakest, and commitment sets that
hour. With commitment fixed they disagree on the hour, although each varies little over the
week. No metric tells more than the committed set of generators, and the
minimum-load snapshot finds the weekly minimum when the fewest units run at minimum load, as under the
largest-first rule.
Under outages, SCR and WSCR follow the change in the network, while ESCR and SDSCR occasionally
report a site stronger while the network weakens. This follows from their interaction term, not
from a convention, although how often it appears in ESCR depends on the construction of the
interaction factor.

Screening over time therefore needs a model of commitment or measured operating states, and an
ESCR value should state how its interaction factor is obtained. Future work should test in EMT
simulation whether the opposing outages are less critical, as ESCR and SDSCR indicate, include converters that inject
reactive current during faults, and extend the study to other networks, IBR penetrations and
plant locations, and a broader set of operating states.
